\chapter{Reference: scroll regions}
\label{ch:refman-scroll}

\begin{aside}
Long content in a small window means a scrollable region.
Keep an offset in the Model; slice the content by the
offset; render only the visible window.
\end{aside}

\section{Working example}

\begin{lstlisting}[language=C++]
#include <maya/maya.hpp>
#include <vector>

using namespace maya;
using namespace maya::dsl;

struct Scroll {
    static constexpr int kVisible = 8;
    struct Model {
        std::vector<std::string> lines;
        int offset = 0;
    };
    struct Up {};
    struct Down {};
    struct PageUp {};
    struct PageDown {};
    struct Quit {};
    using Msg = std::variant<Up, Down, PageUp, PageDown, Quit>;

    static Model init() {
        Model m;
        for (int i = 0; i < 60; ++i) {
            m.lines.push_back(std::format("line {:03d}", i));
        }
        return m;
    }

    static auto update(Model m, Msg msg)
        -> std::pair<Model, Cmd<Msg>>
    {
        int max_off = std::max(0,
            static_cast<int>(m.lines.size()) - kVisible);
        return std::visit(overload{
            [&](Up) {
                m.offset = std::max(0, m.offset - 1);
                return std::pair{m, Cmd<Msg>{}};
            },
            [&](Down) {
                m.offset = std::min(max_off, m.offset + 1);
                return std::pair{m, Cmd<Msg>{}};
            },
            [&](PageUp) {
                m.offset = std::max(0, m.offset - kVisible);
                return std::pair{m, Cmd<Msg>{}};
            },
            [&](PageDown) {
                m.offset = std::min(max_off, m.offset + kVisible);
                return std::pair{m, Cmd<Msg>{}};
            },
            [](Quit) {
                return std::pair{Model{}, Cmd<Msg>::quit()};
            },
        }, msg);
    }

    static Element view(const Model& m) {
        std::vector<Element> visible;
        int start = m.offset;
        int end   = std::min(start + kVisible,
                             static_cast<int>(m.lines.size()));
        for (int i = start; i < end; ++i) {
            visible.push_back(text(m.lines[i]));
        }
        return v(
            v(visible),
            blank_,
            text(std::format("{} / {}",
                m.offset + 1, m.lines.size())) | Dim
        ) | pad<1> | border_<Round>;
    }

    static auto subscribe(const Model&) -> Sub<Msg> {
        return key_map<Msg>({
            {'q', Quit{}},
            {SpecialKey::Up,       Up{}},
            {SpecialKey::Down,     Down{}},
            {SpecialKey::PageUp,   PageUp{}},
            {SpecialKey::PageDown, PageDown{}},
        });
    }
};

int main() { run<Scroll>({.title = "scroll"}); }
\end{lstlisting}

\section{Notes}

\begin{itemize}[leftmargin=*]
  \item Slice in render, not in update; the Model stores
    only the offset.
  \item Clamp the offset both directions; users will mash
    arrow keys past the edges.
  \item For very long content, render a scrollbar from the
    offset and total ratio.
\end{itemize}

\section{Recap}

\begin{recap}
\begin{itemize}[leftmargin=*]
  \item Offset in Model; slice in view.
  \item Clamp at both edges.
  \item Show position as "n / total".
\end{itemize}
\end{recap}

\section*{Workshop}

\begin{enumerate}[leftmargin=*]
  \item Add a scrollbar in the right margin.
  \item Auto-scroll to the bottom when new lines arrive.
  \item Implement \code{g} to jump home, \code{G} to end.
\end{enumerate}
